\documentclass{article}
\usepackage{graphicx} % Required for inserting images

\title{Astra Rendering Engine}
\author{Matthew Chen, Justin Ng, Cindy Liu, Lingnan Meng}
\date{March 2026}

\begin{document}

\maketitle

\textbf{"How can we systematically map continuous 3D mathematical data (f(x,y)→z) into a 2D visual color space using tree-based recursive evaluation?"}

\section{\textbf{Introduction}}

The visualization of multivariable mathematical functions of the form f(x,y)→z is a fundamental tool in mathematics and computer science. However, creating an environment where a user can input an arbitrary string—such as sin(x) + cos(y\^2)—and see it rendered requires significant computational overhead. Python’s built-in evaluation methods are notoriously slow for large-scale computations and present security risks. When rendering a high-resolution 800×800grid, the program must perform hundreds of thousands of evaluations per frame.

    To solve the dual problems of parsing user input and executing calculations efficiently, we treat the user’s input as a hierarchical data structure. By converting the equation into an Abstract Syntax Tree (AST), we represent mathematical operations as nodes and variables/constants as leaves.

\textbf{The goal of this project is to develop a rendering engine that uses an Abstract Syntax Tree to parse user-defined equations, enabling the visualization of complex multivariable functions through recursive tree evaluation.}

\section{\textbf{Datasets}}

Unlike traditional data science projects, our engine does not rely on static CSV or JSON datasets. Instead, our "dataset" is procedurally generated from user input. The input consists of standard mathematical strings (e.g., 43.43*ln(-2x+(1+sqrt(5))/2)). 

Our program generates its own internal dataset: a 2D coordinate grid defined by MATH\_MIN, MATH\_MAX, and GRID\_RESOLUTION. The engine evaluates the AST at every (x,y) coordinate in this grid to produce a matrix of z-values, which are then mapped to RGB color channels.

\section{\textbf{Computational Overview}}

\subsection{\textbf{Data Representation and Trees}}

The central data structure of our project is the Expression Tree, implemented via the Node and Equation classes in the Equation.py file.

\begin{itemize}
    \item \textbf{Internal Nodes:} Represent mathematical operators (e.g., +, -, *, \^, sin, ln).
    \item \textbf{Leaf Nodes:} Represent numerical constants or variables (x, y).
\end{itemize}

\subsection{\textbf{Major Computations}}

\begin{enumerate}
    \item \textbf{Tokenization and Pre-processing:} The tokenize method in our Equation class cleans the raw string, handles LaTeX formatting (like \textbackslash{}frac), and safely isolates variables. Crucially, it algorithms dynamically inject * tokens to handle implicit multiplication (e.g., converting 2x to 2*x) and parses unary negative signs to prevent structural crashes.
    \item \textbf{The Shunting-Yard Algorithm:} We convert standard infix mathematical notation into a binary tree structure based on standard order of operations (PEMDAS). We utilize a stack-based builder (operator\_stack and output\_stack) governed by a PRECEDENCE dictionary to correctly nest operations.
    \item \textbf{Recursive Grid Evaluation:} Inside main.py, a nested loop generates the (x,y) coordinate grid. For each point, the engine calls drawFunc.evaluate(x, y), which triggers a depth-first, post-order traversal of the AST to compute the final z-value.
\end{enumerate}

\subsection{\textbf{Visual and Interactive Reporting}}

We utilize the \textbf{Pygame} library to handle the visual output. The z-values obtained from the AST evaluation are passed to a Color class, which normalizes the data into an integer tuple between (0, 0, 0) and (255, 255, 255). We then use pygame.draw.rect to map these RGB values onto the display surface, creating a continuous heat-map style gradient. 

\section{\textbf{Instructions for Running the Program}}

To run the Astra Render Engine:

\begin{enumerate}
    \item Ensure you have Python 3.13.5 installed.
    \item Install the required external libraries:  install pygame and pyperclip from requirements.txt.
    \item Execute python main.py.
    \item Copy and import Example function from the REQUIREMENTS.txt file into Pygame settings window (instructions below)
\end{enumerate}
\textbf{Expected Output:} Upon running main.py, a Pygame window will open. 
\\\\Extra Notes
\begin{enumerate}
    \item{The Function Tab contains all functions used in the program}
    \item{The Color, Restriction, and Draw tabs all refer to functions in the tab}
\end{enumerate}
NOTE: Higher Resolutions (x points and y points) WILL take longer to render.
\\Copy the Sample Function (ON THE REQUIREMENTS FILE) and click import in Settings. Steps:
\begin{enumerate}
    \item open the Pygame window through main.py
    \item Go to Requirements.txt file and copy text under Example Function to clipboard
    \item go to Pygame Window, and go to settings tab
    \item click import on the bottom left
\end{enumerate}

\section{\textbf{Changes to the Project Plan}}

Our decision to use Pygame came with many challenges. One significant challenge was how to implement textboxes and get inputs into the variables such that we can operate cleanly on it. To do this, an EntryField abstract class was made which allowed us to have a framework for the  Functions, Colors, Restrictions, and Draw subclasses that we would later create in the Main class.

\section{\textbf{Discussion and Results}}
Our primary goal was to build an engine capable of parsing and visualizing complex multivariable functions using ASTs. The results of our computational exploration demonstrate that while Expression Trees are highly effective for respecting mathematical precedence and evaluating raw strings, they introduce significant computational bottlenecks in a single-threaded Python environment.
\\\\Through our testing, we discovered that the time complexity of our rendering engine is O(E⋅N\^2), where E is the total number of nodes in all used ASTs, and N is the resolution of the grid(one side length). Because the grid requires N\^2 evaluations, the tree traversal heavily dominates the runtime. For an 800×800 grid (N=800) and a moderately complex equation set (E=116), our engine performed over 74 million recursive node visits.
\\\\While our engine successfully rendered accurate RGB mappings (verified against the Desmos Graphing Calculator), it took approximately 45 seconds to render a single frame at high resolutions. This limitation is primarily due to Python's memory allocation overhead during deep recursion and the cost of string comparisons at each node. However, this comes at a tradeoff of having better resolution than existing graphing software. To address these limitations in future iterations, we could implement "Constant Folding" (an algorithmic pre-compute pass that collapses static branches of the tree, reducing E) and migrate our Pygame drawing logic to numpy array manipulation to bypass the overhead of individual draw.rect calls. 
\\\\Beyond the computational challenges, our group also encountered significant obstacles in the user interface layer—especially when working with Pygame. Unlike higher-level GUI frameworks, Pygame does not natively provide support for common interface components such as textboxes, buttons, or input fields. This meant that even basic functionality, such as allowing users to type in mathematical expressions or select specific indexes in text fields with a mouse click, had to be implemented entirely from scratch.
One of the primary difficulties was designing a robust textbox system that could handle real-time keyboard input. We needed to manually track focus states (i.e., whether a textbox was “active”), detect mouse clicks within bounding rectangles, and process individual keypress events to update the displayed string. Handling edge cases—such as backspacing, cursor positioning, and managing textoverflow—required additional logic that significantly increased the complexity of our codebase. 
\\\\We initially explored importing external modules to simplify this process, but many were incompatible with pygame or required a different library which conflicted with our initial idea and implementation. As a result, we opted to build a custom textbox implementation tailored specifically to our needs. This involved rendering text surfaces dynamically based on stored scrolled values, recalculating widths, and implementing clipping mechanisms to ensure visual consistency.
\\\\Design considerations also played a role. We needed to balance readability and responsiveness while ensuring that the textbox interface did not interfere with the rendering performance of our engine. This led to iterative refinements in layout, font scaling, and input responsiveness. Overall, while these UI challenges were separate from the core AST engine, they highlighted the importance of integrating user interaction seamlessly with computational systems—an essential step in transforming a functional backend into a usable and interactive application.
\\\\Ultimately, this project successfully proves that compiler-level data structures like ASTs can be bridged with visual graphics libraries to create powerful, customized rendering tools.

\section{\textbf{References}}

\begin{itemize}
    \item GeeksforGeeks. \textit{Expression Tree Data Structure.} \\{https://www.geeksforgeeks.org/dsa/expression-tree/}
    \item Pygame Community. \textit{Pygame draw module documentation.} \\{https://www.pygame.org/docs/ref/draw.html}
    \item Desmos Studio. \textit{Desmos Graphing Calculator} (Used for verification). \\{https://www.desmos.com/}
    \item {https://www.desmos.com/calculator/h3ogiandp1} (basis for test graph)
\end{itemize}

 
\end{document}

\section{Introduction